% Supplementary information (separate PDF). Build: python3 paper/build.py
\documentclass[11pt]{article}
% Shared preamble of manuscript.tex and supplement.tex (built with tectonic; see paper/build.py)
\usepackage[a4paper,margin=2.5cm]{geometry}
\usepackage{fontspec}   % XeTeX (tectonic): Latin Modern OpenType by default, Unicode input
\usepackage{microtype}
\usepackage[running]{lineno}
\usepackage{graphicx}
\usepackage{booktabs}
\usepackage{longtable}
\usepackage{tabularx}
\usepackage{array}
\usepackage{caption}
\captionsetup{font=small,labelfont=bf,labelsep=period}
\usepackage{enumitem}
\setlist{itemsep=1pt,topsep=3pt}
\usepackage{amsmath}
\usepackage{xcolor}
\usepackage[numbers,sort&compress,square]{natbib}
\usepackage[hidelinks]{hyperref}
\urlstyle{same}
\setlength{\parskip}{3pt}
\renewcommand{\arraystretch}{1.15}
\newcolumntype{L}[1]{>{\raggedright\arraybackslash}p{#1}}
\newcolumntype{Y}{>{\raggedright\arraybackslash}X}

% Numbers: every number taken from the data is printed with \nm{key}; the values come from
% numbers.tex, which paper/make_numbers.py writes from the repository's files (sources in
% NUMBERS_generated.tsv, text-to-key map in NUMBERS.md). An unknown key stops the build.
\makeatletter
\newcommand{\nmdef}[2]{\expandafter\def\csname nm@#1\endcsname{#2}}
\newcommand{\nm}[1]{\ifcsname nm@#1\endcsname\csname nm@#1\endcsname\else\errmessage{unknown number key: #1}\fi}
\makeatother
% Generated by paper/make_numbers.py. Do not edit. Sources: paper/NUMBERS_generated.tsv
\nmdef{titin_len}{35,991}
\nmdef{pkz_len}{45,212}
\nmdef{uniprotkb_gt_titin}{51}
\nmdef{uniprotkb_gt_pkz}{4}
\nmdef{uniprotkb_pe3}{37}
\nmdef{uniprotkb_pe4}{14}
\nmdef{uniprotkb_reviewed}{0}
\nmdef{uniparc_gt_titin}{477}
\nmdef{uniparc_gt_pkz}{65}
\nmdef{ncbi_gt_titin}{472}
\nmdef{ncbi_gt_pkz}{47}
\nmdef{refseq_gt_titin}{349}
\nmdef{union_seqs}{552}
\nmdef{union_gt_pkz}{71}
\nmdef{union_live}{523}
\nmdef{excl_nolive}{29}
\nmdef{excl_patent}{6}
\nmdef{excl_partial}{7}
\nmdef{longest_record}{187,500}
\nmdef{n_retained}{510}
\nmdef{n_excluded}{42}
\nmdef{n_retained_gt_pkz}{57}
\nmdef{loci_retained_stage0}{228}
\nmdef{xp_records}{329}
\nmdef{xp_genes}{52}
\nmdef{n_euk}{452}
\nmdef{n_bact}{57}
\nmdef{n_arch}{1}
\nmdef{big_euk}{34}
\nmdef{big_bact}{23}
\nmdef{s1_all_A}{114}
\nmdef{s1_all_B}{213}
\nmdef{s1_all_C}{154}
\nmdef{s1_all_D}{29}
\nmdef{s1_big_A}{0}
\nmdef{s1_big_B}{4}
\nmdef{s1_big_C}{36}
\nmdef{s1_big_D}{17}
\nmdef{s2_all_A}{119}
\nmdef{s2_all_B}{267}
\nmdef{s2_all_C}{95}
\nmdef{s2_all_D}{29}
\nmdef{s2_big_A}{0}
\nmdef{s2_big_B}{5}
\nmdef{s2_big_C}{35}
\nmdef{s2_big_D}{17}
\nmdef{fin_all_A}{120}
\nmdef{fin_all_B}{267}
\nmdef{fin_all_C}{94}
\nmdef{fin_all_D}{29}
\nmdef{fin_big_A}{1}
\nmdef{fin_big_B}{5}
\nmdef{fin_big_C}{34}
\nmdef{fin_big_D}{17}
\nmdef{s1s2_C_to_B}{59}
\nmdef{s1s2_B_to_A}{5}
\nmdef{s1s2_C_to_B_apostichopus}{48}
\nmdef{rule_D1_all}{25}
\nmdef{rule_D1_big}{15}
\nmdef{rule_D2_all}{16}
\nmdef{rule_D2_big}{10}
\nmdef{rule_D3_all}{24}
\nmdef{rule_D3_big}{15}
\nmdef{rule_D5_all}{3}
\nmdef{rule_D5_big}{2}
\nmdef{rule_D1_big_euk}{14}
\nmdef{d1_pattern_all}{72}
\nmdef{d1_pattern_big}{32}
\nmdef{d1_pattern_big_euk}{29}
\nmdef{repeat_flag_all}{47}
\nmdef{repeat_flag_big}{17}
\nmdef{flag_big_weak_coding}{11}
\nmdef{flag_big_weak_coding_range}{0.919--0.942}
\nmdef{flag_big_coding}{4}
\nmdef{flag_big_nodna}{2}
\nmdef{d_big_d1_with_dna}{15}
\nmdef{d_big_d5}{2}
\nmdef{rna_not_tested_d1}{0}
\nmdef{rep_models}{43}
\nmdef{rep_testable}{10}
\nmdef{rep_testable_big}{5}
\nmdef{rep_jobs}{7}
\nmdef{rep_runs}{33}
\nmdef{rep_gb}{135}
\nmdef{rep_schisto1_supported}{32}
\nmdef{rep_schisto1_introns}{45}
\nmdef{rep_schisto1_testable}{32}
\nmdef{rep_schisto1_testable_unsup}{0}
\nmdef{rep_schisto1_unique_reads}{44,113}
\nmdef{rep_schisto1_verdict}{partial}
\nmdef{rep_schisto1_grade}{B}
\nmdef{rep_schisto2_supported}{32}
\nmdef{rep_schisto2_introns}{46}
\nmdef{rep_schisto2_testable}{32}
\nmdef{rep_schisto2_testable_unsup}{0}
\nmdef{rep_schisto2_unique_reads}{44,112}
\nmdef{rep_schisto2_verdict}{partial}
\nmdef{rep_schisto2_grade}{B}
\nmdef{rep_owenia_supported}{9}
\nmdef{rep_owenia_introns}{22}
\nmdef{rep_owenia_testable}{6}
\nmdef{rep_owenia_testable_unsup}{0}
\nmdef{rep_owenia_unique_reads}{673}
\nmdef{rep_owenia_verdict}{weak (unsupported introns lie where reads are too few)}
\nmdef{rep_owenia_grade}{C}
\nmdef{rep_arcto_supported}{1}
\nmdef{rep_arcto_introns}{7}
\nmdef{rep_arcto_testable}{0}
\nmdef{rep_arcto_testable_unsup}{0}
\nmdef{rep_arcto_unique_reads}{85}
\nmdef{rep_arcto_verdict}{weak (unsupported introns lie where reads are too few)}
\nmdef{rep_arcto_grade}{C}
\nmdef{rep_noexpr}{6}
\nmdef{rep_noexpr_max_unique}{34}
\nmdef{rep_contradicted}{0}
\nmdef{rep_to_B}{2}
\nmdef{rep_to_D}{0}
\nmdef{rep_run_results}{33}
\nmdef{rep_run_skips}{0}
\nmdef{rep_schisto_err10016938}{32}
\nmdef{mapp_gas1_uj}{25}
\nmdef{mapp_gas1_n}{2,053}
\nmdef{mapp_gas1_pct}{1}
\nmdef{mapp_gas2_uj}{26}
\nmdef{mapp_gas2_n}{1,273}
\nmdef{mapp_gas2_pct}{2}
\nmdef{mapp_gas3_uj}{27}
\nmdef{mapp_gas3_n}{1,208}
\nmdef{mapp_gas3_pct}{2}
\nmdef{mapp_acn_uj}{48}
\nmdef{mapp_acn_n}{1,035}
\nmdef{mapp_acn_pct}{5}
\nmdef{mapp_syn1_uj}{16}
\nmdef{mapp_syn1_n}{203}
\nmdef{mapp_syn1_pct}{8}
\nmdef{mapp_syn2_uj}{9}
\nmdef{mapp_syn2_n}{134}
\nmdef{mapp_syn2_pct}{7}
\nmdef{mapp_schisto1_uj}{37}
\nmdef{mapp_schisto1_n}{45}
\nmdef{mapp_schisto1_pct}{82}
\nmdef{mapp_owenia_uj}{22}
\nmdef{mapp_owenia_n}{22}
\nmdef{mapp_owenia_pct}{100}
\nmdef{mapp_arcto_uj}{7}
\nmdef{mapp_arcto_n}{7}
\nmdef{mapp_arcto_pct}{100}
\nmdef{ctl_release}{2026\_04}
\nmdef{ctl_n}{176}
\nmdef{ctl_silk_n}{3}
\nmdef{ctl_silk_d1}{1}
\nmdef{ctl_spidroin_n}{4}
\nmdef{ctl_spidroin_d1}{3}
\nmdef{ctl_mucin_n}{15}
\nmdef{ctl_mucin_d1}{0}
\nmdef{ctl_adhesin_n}{13}
\nmdef{ctl_adhesin_d1}{4}
\nmdef{ctl_long_n}{150}
\nmdef{ctl_long_d1}{3}
\nmdef{ctl_d1_flagged}{10}
\nmdef{ctl_d23_flagged}{0}
\nmdef{ctl_flagged_with_cds}{8}
\nmdef{ctl_flagged_fickett_min}{1.094}
\nmdef{ctl_flagged_open_max}{1}
\nmdef{hold_n}{9,428}
\nmdef{hold_trigger}{17}
\nmdef{hold_d1_v1}{17}
\nmdef{hold_d1_cds}{10}
\nmdef{hold_d1_v2}{0}
\nmdef{power_models}{57}
\nmdef{power_reps}{20}
\nmdef{power_shift_0.1_v1_contradicted}{0}
\nmdef{power_shift_0.1_v1_under-supported}{0}
\nmdef{power_shift_0.1_v1_partial}{0.997}
\nmdef{power_shift_0.1_v1_full}{0.003}
\nmdef{power_shift_0.1_v2_contradicted}{0}
\nmdef{power_shift_0.1_v2_under-supported}{0}
\nmdef{power_shift_0.1_v2_partial}{0.997}
\nmdef{power_shift_0.1_v2_full}{0.003}
\nmdef{power_shift_0.2_v1_contradicted}{0}
\nmdef{power_shift_0.2_v1_under-supported}{0}
\nmdef{power_shift_0.2_v1_partial}{1}
\nmdef{power_shift_0.2_v1_full}{0}
\nmdef{power_shift_0.2_v2_contradicted}{0}
\nmdef{power_shift_0.2_v2_under-supported}{0.865}
\nmdef{power_shift_0.2_v2_partial}{0.135}
\nmdef{power_shift_0.2_v2_full}{0}
\nmdef{power_shift_0.3_v1_contradicted}{0}
\nmdef{power_shift_0.3_v1_under-supported}{0}
\nmdef{power_shift_0.3_v1_partial}{1}
\nmdef{power_shift_0.3_v1_full}{0}
\nmdef{power_shift_0.3_v2_contradicted}{0}
\nmdef{power_shift_0.3_v2_under-supported}{1}
\nmdef{power_shift_0.3_v2_partial}{0}
\nmdef{power_shift_0.3_v2_full}{0}
\nmdef{power_shift_0.4_v1_contradicted}{0}
\nmdef{power_shift_0.4_v1_under-supported}{0}
\nmdef{power_shift_0.4_v1_partial}{1}
\nmdef{power_shift_0.4_v1_full}{0}
\nmdef{power_shift_0.4_v2_contradicted}{0}
\nmdef{power_shift_0.4_v2_under-supported}{1}
\nmdef{power_shift_0.4_v2_partial}{0}
\nmdef{power_shift_0.4_v2_full}{0}
\nmdef{power_shift_0.5_v1_contradicted}{0.725}
\nmdef{power_shift_0.5_v1_under-supported}{0.198}
\nmdef{power_shift_0.5_v1_partial}{0.076}
\nmdef{power_shift_0.5_v1_full}{0}
\nmdef{power_shift_0.5_v2_contradicted}{0.725}
\nmdef{power_shift_0.5_v2_under-supported}{0.275}
\nmdef{power_shift_0.5_v2_partial}{0}
\nmdef{power_shift_0.5_v2_full}{0}
\nmdef{power_shift_0.6_v1_contradicted}{1}
\nmdef{power_shift_0.6_v1_under-supported}{0}
\nmdef{power_shift_0.6_v1_partial}{0}
\nmdef{power_shift_0.6_v1_full}{0}
\nmdef{power_shift_0.6_v2_contradicted}{1}
\nmdef{power_shift_0.6_v2_under-supported}{0}
\nmdef{power_shift_0.6_v2_partial}{0}
\nmdef{power_shift_0.6_v2_full}{0}
\nmdef{power_skip_0.2_v1_contradicted}{0}
\nmdef{power_skip_0.2_v1_under-supported}{0}
\nmdef{power_skip_0.2_v1_partial}{1}
\nmdef{power_skip_0.2_v1_full}{0}
\nmdef{power_skip_0.2_v2_contradicted}{0}
\nmdef{power_skip_0.2_v2_under-supported}{0.278}
\nmdef{power_skip_0.2_v2_partial}{0.722}
\nmdef{power_skip_0.2_v2_full}{0}
\nmdef{power_skip_0.3_v1_contradicted}{0}
\nmdef{power_skip_0.3_v1_under-supported}{0}
\nmdef{power_skip_0.3_v1_partial}{1}
\nmdef{power_skip_0.3_v1_full}{0}
\nmdef{power_skip_0.3_v2_contradicted}{0}
\nmdef{power_skip_0.3_v2_under-supported}{1}
\nmdef{power_skip_0.3_v2_partial}{0}
\nmdef{power_skip_0.3_v2_full}{0}
\nmdef{power_skip_0.5_v1_contradicted}{0.242}
\nmdef{power_skip_0.5_v1_under-supported}{0.061}
\nmdef{power_skip_0.5_v1_partial}{0.696}
\nmdef{power_skip_0.5_v1_full}{0}
\nmdef{power_skip_0.5_v2_contradicted}{0.242}
\nmdef{power_skip_0.5_v2_under-supported}{0.758}
\nmdef{power_skip_0.5_v2_partial}{0}
\nmdef{power_skip_0.5_v2_full}{0}
\nmdef{power_skip_0.6_v1_contradicted}{0.981}
\nmdef{power_skip_0.6_v1_under-supported}{0.009}
\nmdef{power_skip_0.6_v1_partial}{0.011}
\nmdef{power_skip_0.6_v1_full}{0}
\nmdef{power_skip_0.6_v2_contradicted}{0.981}
\nmdef{power_skip_0.6_v2_under-supported}{0.019}
\nmdef{power_skip_0.6_v2_partial}{0}
\nmdef{power_skip_0.6_v2_full}{0}
\nmdef{pw_contr_ge60_min_pct}{98}
\nmdef{pw_shift50_pct}{73}
\nmdef{pw_skip50_pct}{24}
\nmdef{pw_v1_partial_2040_min_pct}{99.9}
\nmdef{pw_v2_down_3040_min_pct}{100}
\nmdef{pw_v2_shift20_pct}{87}
\nmdef{pw_v2_skip20_pct}{28}
\nmdef{pw_v2_10_down_pct}{0}
\nmdef{dnactl_v3_a_deletion_ill}{against: 1 single-copy boundary(ies) with 0 spanning pairs between well-covered windows; 1 clip pile-up(s) where physical coverage drops; discordant pairs up to 19\% in a window [0.66 kb in 1 stretch(es) where this library cannot place pairs uniquely (repeat), not crossed]}
\nmdef{dnactl_v3_a_deletion_illont}{against: 1 single-copy boundary(ies) with 0 spanning pairs between well-covered windows; 1 clip pile-up(s) where physical coverage drops; discordant pairs up to 19\% in a window [0.66 kb in 1 stretch(es) where this library cannot place pairs uniquely (repeat), not crossed]}
\nmdef{dnactl_v3_b_collapse_ill}{supported by the reads}
\nmdef{dnactl_v3_b_collapse_illont}{supported by the reads}
\nmdef{dnactl_v3_c_expansion_ill}{supported outside 9.37 kb of repeat (no long reads to resolve it)}
\nmdef{dnactl_v3_c_expansion_illont}{supported outside 9.37 kb of repeat; long reads do not span it}
\nmdef{dnactl_v3_d_rrna_correct_ill}{supported outside 4.89 kb of repeat (no long reads to resolve it)}
\nmdef{dnactl_v3_d_rrna_correct_illont}{supported by the reads (short reads; the 4.89 kb repeat block(s) spanned by 13+ long reads)}
\nmdef{dnactl_v3_e_segdup_correct_ill}{inconclusive: only 98.7\% of single-copy gene bases at depth $\geq$ 5 [12.0 kb in multi-copy block(s) that short reads cannot place; 13.04 kb in 4 stretch(es) where this library cannot place pairs uniquely (repeat), not crossed]}
\nmdef{dnactl_v3_e_segdup_correct_illont}{inconclusive: only 98.7\% of single-copy gene bases at depth $\geq$ 5 [12.0 kb in multi-copy block(s) that short reads cannot place; 13.04 kb in 4 stretch(es) where this library cannot place pairs uniquely (repeat), not crossed]}
\nmdef{dnactl_v31_a_deletion_ill}{against: 1 single-copy boundary(ies) with 0 spanning pairs between well-covered windows; 1 clip pile-up(s) where physical coverage drops; discordant pairs up to 19\% in a window [0.66 kb in 1 stretch(es) where this library cannot place pairs uniquely (repeat), not crossed]}
\nmdef{dnactl_v31_a_deletion_illont}{against: 1 single-copy boundary(ies) with 0 spanning pairs between well-covered windows; 1 clip pile-up(s) where physical coverage drops; discordant pairs up to 19\% in a window [0.66 kb in 1 stretch(es) where this library cannot place pairs uniquely (repeat), not crossed]}
\nmdef{dnactl_v31_b_collapse_ill}{against: copy-number anomaly (v3.1): 5 kb at 1.98x, excess copies (collapsed repeat?)}
\nmdef{dnactl_v31_b_collapse_ill_ratio}{1.98}
\nmdef{dnactl_v31_b_collapse_illont}{against: copy-number anomaly (v3.1): 5 kb at 1.98x, excess copies (collapsed repeat?)}
\nmdef{dnactl_v31_b_collapse_illont_ratio}{1.98}
\nmdef{dnactl_v31_c_expansion_ill}{against: copy-number anomaly (v3.1): 10 kb at 0.49x, deficit (expanded or absent?) [8.0 kb in multi-copy block(s) that short reads cannot place; 8.74 kb in 2 stretch(es) where this library cannot place pairs uniquely (repeat), not crossed]}
\nmdef{dnactl_v31_c_expansion_ill_ratio}{0.49}
\nmdef{dnactl_v31_c_expansion_illont}{against: copy-number anomaly (v3.1): 10 kb at 0.49x, deficit (expanded or absent?) [8.0 kb in multi-copy block(s) that short reads cannot place; 8.74 kb in 2 stretch(es) where this library cannot place pairs uniquely (repeat), not crossed]}
\nmdef{dnactl_v31_c_expansion_illont_ratio}{0.49}
\nmdef{dnactl_v31_d_rrna_correct_ill}{supported outside 4.89 kb of repeat (no long reads to resolve it)}
\nmdef{dnactl_v31_d_rrna_correct_illont}{supported by the reads (short reads; the 4.89 kb repeat block(s) spanned by 13+ long reads)}
\nmdef{dnactl_v31_e_segdup_correct_ill}{inconclusive: only 98.7\% of single-copy gene bases at depth $\geq$ 5 [12.0 kb in multi-copy block(s) that short reads cannot place; 13.04 kb in 4 stretch(es) where this library cannot place pairs uniquely (repeat), not crossed]}
\nmdef{dnactl_v31_e_segdup_correct_illont}{inconclusive: only 98.7\% of single-copy gene bases at depth $\geq$ 5 [12.0 kb in multi-copy block(s) that short reads cannot place; 13.04 kb in 4 stretch(es) where this library cannot place pairs uniquely (repeat), not crossed]}
\nmdef{dnactl_real_unchanged}{yes}
\nmdef{dnactl_repeat_depth_ratio_range}{0.84--0.97}
\nmdef{dnactl_omni_repeat_depth_ratio_range}{0.86--0.97}
\nmdef{sens_variants}{33}
\nmdef{sens_big_ab_stable}{31}
\nmdef{sens_big_ab_changed}{D3 open frames $\geq$ 1 (published 2), spliced reads per intron $\geq$ 10 (published 3)}
\nmdef{sens_s6_lenient_all_B}{267}
\nmdef{sens_base_all_B}{267}
\nmdef{sens_d1v1_big_D}{34}
\nmdef{sens_d1v1_all_D}{76}
\nmdef{sens_allA_range}{116--178}
\nmdef{dna_sens_settings}{7}
\nmdef{dna_sens_grade_changes}{min\_pairs\_per\_boundary=10}
\nmdef{prot_spectra}{594,284}
\nmdef{prot_db_entries}{151,457}
\nmdef{prot_fractions}{60}
\nmdef{prot_scan_share_pct}{48}
\nmdef{prot_id_rate_pct}{84.6}
\nmdef{prot_distal}{8,279, 12,982, 15,831, 16,949}
\nmdef{prot_psms}{502,718}
\nmdef{prot_decoy_psms}{5,026}
\nmdef{prot_peptides}{42,000}
\nmdef{prot_proteins}{7,981}
\nmdef{prot_giant_len}{39,686}
\nmdef{prot_giant_peptides}{186}
\nmdef{prot_giant_shared}{0}
\nmdef{prot_giant_psms}{1,487}
\nmdef{prot_giant_residues}{2,357}
\nmdef{prot_giant_cov_pct}{5.9}
\nmdef{prot_giant_regions}{86}
\nmdef{prot_giant_first}{292}
\nmdef{prot_giant_last_end}{16,957}
\nmdef{prot_giant_last_pct}{42.7}
\nmdef{prot_giant_fifths}{182, 3, 1, 0, 0}
\nmdef{prot_giant_first6000}{182}
\nmdef{prot_giant_after6000}{8,279, 12,982, 15,831, 16,949}
\nmdef{prot_giant_protein_q}{0.0002}
\nmdef{prot_rev_q}{0.063}
\nmdef{prot_shuf_q}{0.17}
\nmdef{prot_rev_psms}{0}
\nmdef{prot_shuf_psms}{0}
\nmdef{prot_tryptic_range}{41--63}
\nmdef{prot_repeat_span}{17,653--25,339}
\nmdef{prot_repeat_blocks}{5}
\nmdef{prot_peptides_in_or_after_repeats}{0}
\nmdef{prot_eftu_peptides}{23}
\nmdef{prot_eftu_cov_pct}{55}
\nmdef{prot_eftu_len}{400}
\nmdef{prot_l7_peptides}{6}
\nmdef{prot_l7_cov_pct}{68}
\nmdef{prot_l7_len}{128}
\nmdef{prot_top_n}{12}
\nmdef{prot_top_rank_giant}{11}
\nmdef{prot_top_cov_range}{40.8--78.3}
\nmdef{prot_wall_min}{35}
\nmdef{prot_rss_gb}{5.2}
\nmdef{raw_files}{60}
\nmdef{raw_conversion_errors}{0}
\nmdef{raw_ms2}{1,106,517}
\nmdef{raw_mgf_share_pct}{54}
\nmdef{raw_instrument}{Orbitrap Velos Pro}
\nmdef{raw_psms}{507,612}
\nmdef{raw_peptides}{40,845}
\nmdef{raw_proteins}{7,560}
\nmdef{raw_giant_peptides}{182}
\nmdef{raw_giant_psms}{1,505}
\nmdef{raw_giant_residues}{2,282}
\nmdef{raw_giant_cov_pct}{5.8}
\nmdef{raw_giant_fifths}{179, 3, 0, 0, 0}
\nmdef{raw_giant_first6000}{179}
\nmdef{raw_giant_after6000}{8,279, 12,982, 15,831}
\nmdef{raw_giant_last_end}{15,840}
\nmdef{raw_giant_last_pct}{39.9}
\nmdef{raw_giant_protein_q}{0.00022}
\nmdef{raw_rev_q}{0.089}
\nmdef{raw_shuf_q}{0.17}
\nmdef{raw_rev_psms}{0}
\nmdef{raw_shuf_psms}{0}
\nmdef{raw_eftu_peptides}{24}
\nmdef{raw_eftu_cov_pct}{57}
\nmdef{raw_l7_peptides}{5}
\nmdef{raw_l7_cov_pct}{56}
\nmdef{raw_both}{179}
\nmdef{raw_only_raw}{3}
\nmdef{raw_only_mgf}{7}
\nmdef{raw_mgf_only_after6000}{16,949}
\nmdef{raw_in_or_after_repeats}{0}
\nmdef{raw_wall_min}{56}
\nmdef{prot_peak_last_q}{0.0096}
\nmdef{sy_ms2}{131,021}
\nmdef{sy_raw_files}{5}
\nmdef{sy_db_kept}{99,767}
\nmdef{sy_db_authors}{99,773}
\nmdef{sy_fragments}{6}
\nmdef{sy_fragment_max_len}{46,359}
\nmdef{sy_psms_spec}{272}
\nmdef{sy_peptides_spec}{88}
\nmdef{sy_min_pep_q}{0.018}
\nmdef{sy_min_prot_q}{0.024}
\nmdef{sy_authors_n}{37}
\nmdef{sy_authors_found}{31}
\nmdef{sy_giant_psms_spec}{0}
\nmdef{sy_giant_psms_any}{17}
\nmdef{sy_rev_psms_any}{27}
\nmdef{sy_shuf_psms_any}{15}
\nmdef{sy_calc3_psms}{30}
\nmdef{sy_tub_psms}{37}
\nmdef{sy_ef2_psms}{9}
\nmdef{prot_authors_file}{Mudpit\_160213\_OV3\_P5\_SK\_JHa\_372-Pe1\_01.mzid.gz}
\nmdef{prot_authors_len}{39,671}
\nmdef{prot_authors_peptides}{174}
\nmdef{prot_authors_fifths}{130, 2, 2, 2, 0}
\nmdef{prot_authors_last_start}{25,870}
\nmdef{prot_authors_pass_peptides}{130}
\nmdef{prot_authors_pass_fifths}{101, 0, 0, 0, 0}
\nmdef{pride_omni}{0}
\nmdef{pride_noc881}{0}
\nmdef{pride_noc1730}{0}
\nmdef{pride_sycon}{2}
\nmdef{pride_velam}{1}
\nmdef{prot_raw_files}{60}
\nmdef{prot_raw_gb}{15.3}
\nmdef{prot_closed_genome_proteins}{1,863}
\nmdef{raw_giant_last_pct_round}{40}
\nmdef{rep_untestable_noruns}{28}
\nmdef{rep_untestable_nolocus}{5}
\nmdef{twin_fix_n}{6}
\nmdef{twin_fix_big}{0}
\nmdef{rep_noruns_plectranthus}{25}
\nmdef{rep_noruns_merlangius}{2}
\nmdef{rep_noruns_coilia}{1}
\nmdef{raw_giant_unc_pct}{60}
\nmdef{prot_giant_grade}{C}
\nmdef{lb_loci}{213}
\nmdef{lb_excluded_rows}{42}
\nmdef{lb_A}{9}
\nmdef{lb_B}{85}
\nmdef{lb_C}{90}
\nmdef{lb_D}{29}
\nmdef{lb_big}{55}
\nmdef{lb_big_A}{1}
\nmdef{lb_big_B}{3}
\nmdef{lb_big_C}{34}
\nmdef{lb_big_D}{17}
\nmdef{lb_merged_records}{11}
\nmdef{lb_A_titin_loci}{8}
\nmdef{lb_A_titin_min}{36,247}
\nmdef{lb_A_titin_max}{38,490}
\nmdef{rna_runs_total}{269}
\nmdef{rna_runs_kept}{261}
\nmdef{rna_runs_superseded}{8}
\nmdef{rna_runs_followup}{16}
\nmdef{rna_species}{71}
\nmdef{rna_runs_v1}{8}
\nmdef{rna_runs_v2}{253}
\nmdef{rna_reads_billion}{11.6}
\nmdef{rna_runs_stream_error}{51}
\nmdef{rna_runs_stream_error_first}{46}
\nmdef{rna_stream_error_frac_median_pct}{68}
\nmdef{rna_stream_error_frac_min_pct}{9}
\nmdef{rna_runs_saturation}{55}
\nmdef{rna_runs_readcap}{20}
\nmdef{rna_frac_read_median_pct}{88}
\nmdef{rna_runs_lt50}{49}
\nmdef{rna_full}{59}
\nmdef{rna_partial}{75}
\nmdef{rna_weak}{27}
\nmdef{rna_noexpr}{12}
\nmdef{rna_noruns}{127}
\nmdef{rna_nottested}{15}
\nmdef{rna_notselected}{137}
\nmdef{rna_prok}{58}
\nmdef{rna_contradicted}{0}
\nmdef{rna_tested_models}{173}
\nmdef{titin_runs}{6}
\nmdef{titin_reads_unique}{361,467}
\nmdef{titin_n_introns}{361}
\nmdef{titin_introns_unique_support}{340}
\nmdef{titin_introns_multi_only}{13}
\nmdef{titin_introns_counted_supported}{353}
\nmdef{titin_cds_breadth}{0.973}
\nmdef{titin_testable_introns}{325}
\nmdef{titin_testable_introns_unsupported}{0}
\nmdef{pkz_runs}{7}
\nmdef{pkz_reads_unique}{71,183}
\nmdef{pkz_n_introns}{16}
\nmdef{pkz_introns_unique_support}{9}
\nmdef{pkz_introns_multi_only}{7}
\nmdef{pkz_introns_counted_supported}{16}
\nmdef{pkz_cds_breadth}{0.733}
\nmdef{pkz_testable_introns}{12}
\nmdef{pkz_testable_introns_unsupported}{0}
\nmdef{sycon_runs}{6}
\nmdef{sycon_reads_unique}{49,375}
\nmdef{sycon_n_introns}{24}
\nmdef{sycon_introns_unique_support}{22}
\nmdef{sycon_introns_multi_only}{0}
\nmdef{sycon_introns_counted_supported}{22}
\nmdef{sycon_cds_breadth}{1.0}
\nmdef{sycon_testable_introns}{24}
\nmdef{sycon_testable_introns_unsupported}{2}
\nmdef{halichondria_runs}{4}
\nmdef{halichondria_reads_unique}{70,304}
\nmdef{halichondria_n_introns}{2}
\nmdef{halichondria_introns_unique_support}{2}
\nmdef{halichondria_introns_multi_only}{0}
\nmdef{halichondria_introns_counted_supported}{2}
\nmdef{halichondria_cds_breadth}{0.979}
\nmdef{halichondria_testable_introns}{0}
\nmdef{halichondria_testable_introns_unsupported}{0}
\nmdef{titin_introns_none}{8}
\nmdef{microcaecilia_testable}{41}
\nmdef{microcaecilia_testable_unsup}{1}
\nmdef{microcaecilia_verdict}{partial}
\nmdef{rana_testable}{63}
\nmdef{rana_testable_unsup}{1}
\nmdef{rana_verdict}{partial}
\nmdef{solea_testable}{57}
\nmdef{solea_testable_unsup}{1}
\nmdef{solea_verdict}{partial}
\nmdef{pantherophis_testable}{21}
\nmdef{pantherophis_testable_unsup}{0}
\nmdef{pantherophis_verdict}{weak (unsupported introns lie where reads are too few)}
\nmdef{liftover_moved}{143}
\nmdef{liftover_n_introns}{334}
\nmdef{cam_after_sup}{67}
\nmdef{cam_after_testable_unsup}{0}
\nmdef{ovis_length}{98,182}
\nmdef{ovis_n_exons}{37}
\nmdef{ovis_genomic_span_bp}{356,419}
\nmdef{ovis_protein_period_aa}{22}
\nmdef{ovis_protein_period_identity}{0.847}
\nmdef{ovis_rep25_frac}{0.9922}
\nmdef{ovis_gc_cds}{0.678}
\nmdef{ovis_dna_tandem_lag_bp}{22}
\nmdef{ovis_dna_tandem_identity}{0.939}
\nmdef{ovis_fickett_median}{0.465}
\nmdef{ovis_n_other_frames_open}{2}
\nmdef{ovis_min_stops_per_kcodon_other_frames}{1.2}
\nmdef{ovis_rna_reads}{50,777,537}
\nmdef{ovis_rna_candidates}{0}
\nmdef{ovis_period22_entries}{10}
\nmdef{cirrhinus_len}{59,714}
\nmdef{cirrhinus_exons}{454}
\nmdef{cirrhinus_runs}{6}
\nmdef{cirrhinus_reads_unique}{54,161}
\nmdef{cirrhinus_n_introns}{453}
\nmdef{cirrhinus_introns_counted_supported}{356}
\nmdef{bagarius_len}{54,498}
\nmdef{bagarius_exons}{417}
\nmdef{bagarius_runs}{2}
\nmdef{bagarius_reads_unique}{4,277}
\nmdef{bagarius_n_introns}{416}
\nmdef{bagarius_introns_counted_supported}{144}
\nmdef{cirrhinus_second_copy_start_aa}{31,001}
\nmdef{cirrhinus_first_ttna}{30 of 30}
\nmdef{cirrhinus_second_ttnb}{20 of 29}
\nmdef{bagarius_second_copy_start_aa}{32,001}
\nmdef{bagarius_first_ttna}{31 of 31}
\nmdef{bagarius_second_ttnb}{11 of 22}
\nmdef{cirrhinus_join_intron_aa}{30,675}
\nmdef{cirrhinus_join_intron_reads}{0}
\nmdef{cirrhinus_join_before_median}{103}
\nmdef{cirrhinus_join_after_median}{2}
\nmdef{ovis_entries}{11}
\nmdef{ovis_period_multiple22}{11}
\nmdef{sycon_gap_patch_bp}{2,968}
\nmdef{sycon_best_homolog_len}{4,063}
\nmdef{sycon_len}{50,516}
\nmdef{nbc00881_len}{59,275}
\nmdef{nbc00881_ratio}{1.27}
\nmdef{nbc00881_best_len}{46,734}
\nmdef{nbc00881_orf_bp}{177,828}
\nmdef{nbc01730_len}{51,778}
\nmdef{nbc01730_ratio}{1.11}
\nmdef{nbc01730_best_len}{46,734}
\nmdef{nbc01730_orf_bp}{155,337}
\nmdef{pkz_uniprot_pe}{Predicted}
\nmdef{stage1_cpu_hours}{4.4}
\nmdef{ovis_stops_six_frames}{0, 1.2, 1.2, 45.6, 45.2, 45.5}
\nmdef{dna_oub_reads}{83,321,080}
\nmdef{dna_oub_reads_expected}{83,321,080}
\nmdef{dna_oub_depth}{45}
\nmdef{dna_oub_ratio}{1.12}
\nmdef{dna_oub_span_median}{38}
\nmdef{dna_oub_span_min_testable}{4}
\nmdef{dna_oub_discordant_pct}{1.5}
\nmdef{dna_oub_insert}{288}
\nmdef{dna_oub_untestable}{1}
\nmdef{dna_oub_boundaries}{257}
\nmdef{dna_oub_verdict}{inconclusive: weakest single-copy boundary spanned by 4 [1.34 kb in 2 stretch(es) where this library cannot place pairs uniquely (repeat), not crossed]}
\nmdef{dna_oub_control}{yes}
\nmdef{dna_oub_v2}{against: 1 single-copy boundary(ies) with 0 spanning pairs between well-covered windows}
\nmdef{dna_oub_E_at_untestable}{0}
\nmdef{dna_oub_obs_at_untestable}{0}
\nmdef{dna_oub_weakest}{4--5.8}
\nmdef{dna_nbc00881_reads}{14,056,286}
\nmdef{dna_nbc00881_reads_expected}{14,056,286}
\nmdef{dna_nbc00881_depth}{252}
\nmdef{dna_nbc00881_ratio}{1.02}
\nmdef{dna_nbc00881_span_median}{252}
\nmdef{dna_nbc00881_span_min_testable}{159}
\nmdef{dna_nbc00881_discordant_pct}{1.2}
\nmdef{dna_nbc00881_insert}{340}
\nmdef{dna_nbc00881_untestable}{1}
\nmdef{dna_nbc00881_boundaries}{177}
\nmdef{dna_nbc00881_verdict}{supported by the reads (short reads; the 1.1 kb repeat block(s) spanned by 64+ long reads)}
\nmdef{dna_nbc00881_control}{yes}
\nmdef{dna_nbc00881_v2}{supported by the reads}
\nmdef{dna_nbc00881_E_at_untestable}{0}
\nmdef{dna_nbc00881_obs_at_untestable}{6}
\nmdef{dna_nbc00881_weakest}{159--163}
\nmdef{dna_nbc00881_ont_reads}{61,843}
\nmdef{dna_nbc00881_ont_span_blocks}{64--67}
\nmdef{dna_nbc01730_reads}{15,565,314}
\nmdef{dna_nbc01730_reads_expected}{15,565,314}
\nmdef{dna_nbc01730_depth}{272}
\nmdef{dna_nbc01730_ratio}{1.02}
\nmdef{dna_nbc01730_span_median}{284}
\nmdef{dna_nbc01730_span_min_testable}{74}
\nmdef{dna_nbc01730_discordant_pct}{1.4}
\nmdef{dna_nbc01730_insert}{347}
\nmdef{dna_nbc01730_untestable}{8}
\nmdef{dna_nbc01730_boundaries}{155}
\nmdef{dna_nbc01730_verdict}{supported by the reads (short reads; the 9.0 kb repeat block(s) spanned by 9+ long reads)}
\nmdef{dna_nbc01730_control}{yes}
\nmdef{dna_nbc01730_v2}{supported by the reads (short reads; the 9.0 kb repeat block(s) spanned by 9+ long reads)}
\nmdef{dna_nbc01730_E_at_untestable}{0, 0, 0, 0, 0, 0, 0, 0}
\nmdef{dna_nbc01730_obs_at_untestable}{0, 0, 0, 0, 0, 0, 0, 0}
\nmdef{dna_nbc01730_weakest}{74--22.2}
\nmdef{dna_nbc01730_ont_reads}{36,579}
\nmdef{dna_nbc01730_ont_span_blocks}{9}
\nmdef{dna_tdt_reads}{85,214,796}
\nmdef{dna_tdt_reads_expected}{85,214,796}
\nmdef{dna_tdt_depth}{40}
\nmdef{dna_tdt_ratio}{1.08}
\nmdef{dna_tdt_span_median}{34}
\nmdef{dna_tdt_span_min_testable}{14}
\nmdef{dna_tdt_discordant_pct}{2.1}
\nmdef{dna_tdt_insert}{295}
\nmdef{dna_tdt_untestable}{1}
\nmdef{dna_tdt_boundaries}{257}
\nmdef{dna_tdt_verdict}{supported outside 1.32 kb of repeat (no long reads to resolve it)}
\nmdef{dna_tdt_control}{yes}
\nmdef{dna_tdt_v2}{n/a (judged under v3 only)}
\nmdef{dna_tdt_E_at_untestable}{0}
\nmdef{dna_tdt_obs_at_untestable}{0}
\nmdef{dna_tdt_weakest}{14--19.1}
\nmdef{dna_pmg_reads}{78,636,086}
\nmdef{dna_pmg_reads_expected}{78,636,086}
\nmdef{dna_pmg_depth}{52}
\nmdef{dna_pmg_ratio}{1.06}
\nmdef{dna_pmg_span_median}{45}
\nmdef{dna_pmg_span_min_testable}{15}
\nmdef{dna_pmg_discordant_pct}{1.9}
\nmdef{dna_pmg_insert}{292}
\nmdef{dna_pmg_untestable}{1}
\nmdef{dna_pmg_boundaries}{257}
\nmdef{dna_pmg_verdict}{supported outside 1.3 kb of repeat (no long reads to resolve it)}
\nmdef{dna_pmg_control}{yes}
\nmdef{dna_pmg_v2}{against: 1 single-copy boundary(ies) with 0 spanning pairs between well-covered windows}
\nmdef{dna_pmg_E_at_untestable}{0}
\nmdef{dna_pmg_obs_at_untestable}{0}
\nmdef{dna_pmg_weakest}{15--43.6}
\nmdef{orf_bp}{257,412}
\nmdef{contig_len_tdt}{1,614,705}
\nmdef{rep_lag_bp}{10,362}
\nmdef{rep_lag_codons}{3,454}
\nmdef{rep_copy_bp}{1,184}
\nmdef{rep_copy_a}{223,454--224,637}
\nmdef{rep_copy_b}{233,816--234,999}
\nmdef{rep_diffs}{2}
\nmdef{orf_len_aa}{85,803}
\nmdef{rep_copy_bp_kb}{1.18}
\nmdef{dna_oub_weakest_obs}{4}
\nmdef{dna_oub_weakest_exp}{6}
\nmdef{orf_len_if_excised}{82,349}
\nmdef{omni_copies}{3}
\nmdef{omni_aa_identity_min_pct}{99.96}
\nmdef{omni_aa_identity_max_pct}{99.98}
\nmdef{omni_contig_lengths}{1,614,705, 1,613,271, 1,168,148}
\nmdef{omni_n_in_cds}{0}
\nmdef{hmm_panel}{74}
\nmdef{s2_A_titin}{119}
\nmdef{s1_fix_B_to_C}{48}
\nmdef{titin_downgrades_withdrawn}{5}
\nmdef{cirrhinus_first_end}{31,000}


\newcommand{\gene}[1]{\textit{#1}}
\newcommand{\spn}[1]{\textit{#1}}
\newcommand{\Ca}{\textit{Candidatus}}
\newcommand{\placeholder}[1]{\textcolor{red!70!black}{[#1]}}
% table bodies from paper/tables are read with the primitive \input (LaTeX's \input breaks \bottomrule)
\makeatletter\newcommand{\tablebody}[1]{\@@input #1 }\makeatother
\newcommand{\file}[1]{\texttt{\detokenize{#1}}}

\renewcommand{\thetable}{S\arabic{table}}
\renewcommand{\thefigure}{S\arabic{figure}}
\usepackage{placeins}
\setcounter{table}{1}  % Table S1 is the leaderboard file; the first table below is S2
\title{Supplementary information for ``Evidence grades for protein database entries longer than titin''}
\author{Ahmet Tahir Mat\\[2pt]\small Papex Labs}
\date{\small Version 2, 18 November 2026 (version 1: 13 November 2026). Preprint, not peer reviewed.}

\begin{document}
\maketitle
\linenumbers
\tableofcontents

\section*{Supplementary Note 1. Rule revisions and additions}
\addcontentsline{toc}{section}{Note 1. Rule revisions and additions}

The rules were set by us on two controls (human titin NP\_001254479.2 and PKZILLA-1) and revised seven times while the
work was in progress. Each revision, with its effect:

\begin{enumerate}[leftmargin=*]
\item \textbf{Isoforms counted as independent support.} The ``other assembly'' line (S2) first counted isoforms of one gene
and same-assembly entries as independent assemblies. Fixing this moved \nm{s1_fix_B_to_C} sequences from B to C at Stage 1.
\item \textbf{Gene-wide junction test and saturation stop.} A gene-wide test of missing junction support, combined with a
read-saturation stop that ended runs early when coverage was concentrated at the 3$'$ end, first downgraded
\nm{titin_downgrades_withdrawn} titin models from B to C. We withdrew the test, now count missing support only at introns
whose flanking exons are covered, and re-ran the affected runs without the stop.
\item \textbf{The 1-kb window repeat mask.} In the source-read check this mask produced two false ``against'' verdicts for
the Omnitrophota ORF (Note~3). It was replaced by a library simulation without changing any threshold.
\item \textbf{The repeat rule.} Positive controls showed that the repeat pattern alone also flags real silk proteins and
adhesins (Note~2). It now counts only with DNA-level evidence (D2 or D3). This moved \nm{repeat_flag_all} sequences from D
to C, \nm{repeat_flag_big} of them above \nm{pkz_len} aa.
\item \textbf{The RNA-seq verdict order.} A power test showed that models with 20--40\,\% wrong introns were called
``partial''. The test for unsupported covered introns now runs first. No real verdict changed.
\item \textbf{A copy-number test in the source-read check (v3.1).} Known misassemblies showed that a collapsed or duplicated
repeat passed the check. The added test catches both and changed no real verdict.
\item \textbf{Twin records.} The orthologue and other-assembly lines (S1, S2) counted a RefSeq record and its GenBank or
UniProtKB twin from the same genome as independent; for S1 this happened when one species name carried ``Candidatus'' and
the other did not. Fixing this moved \nm{twin_fix_n} sequences from B to C, \nm{twin_fix_big} of them longer than \nm{pkz_len} aa. The error was found during
the proteomics analysis (WP\_128698900.1 had been B on an ``orthologue'' that was the UniProtKB copy of its own genome).
\end{enumerate}

The repeat rule, the RNA verdict order and the copy-number test were revised after we saw the controls. Only the repeat rule
was tested again on an independent set (the \nm{hold_n}-protein hold-out set); the other two have not been tested on fresh
cases.

Two rules were added late. \textbf{The peptide line (S6)} was written after we had seen its only test case, the search of the
deposited PXD025008 peak lists. A first draft asked for unique peptides in three of the five length fifths, which
WP\_128698900.1 met in that search with \nm{prot_giant_first6000} of its \nm{prot_giant_peptides} peptides in the first
6,000 aa. We replaced it with all five fifths, because three would count N-terminal peptides as support for the whole
model. Under the first draft, that entry and its GenBank twin would have been B; in the later search of the raw files,
peptides cover only two fifths, so neither version gives B. \textbf{The guard for the repeat-flag RNA-seq stage} (Note~4)
was written and implemented before that stage's data existed.

\section*{Supplementary Note 2. Controls of the grading rules}
\addcontentsline{toc}{section}{Note 2. Controls of the grading rules}

We added four controls after the audit had run, to test the rules rather than the entries. Three of them changed a rule.

\paragraph{Real repetitive proteins and the repeat rule.} Before scoring, we fixed a design set of real proteins with
protein-level evidence in UniProtKB (release \nm{ctl_release}), \nm{ctl_n} entries in all: \nm{ctl_silk_n} reviewed silk
proteins and \nm{ctl_spidroin_n} complete spidroins, \nm{ctl_mucin_n} mucins, \nm{ctl_adhesin_n} bacterial adhesins and the
\nm{ctl_long_n} reviewed proteins of at least 5,000 aa. The pattern alone flagged \nm{ctl_d1_flagged}:
\nm{ctl_spidroin_d1} of the \nm{ctl_spidroin_n} complete spidroins, the \spn{Bombyx mori} fibroin heavy chain,
\nm{ctl_adhesin_d1} of the \nm{ctl_adhesin_n} adhesins (all serine-rich repeat proteins) and two long human proteins, the
IgGFc-binding protein and NBPF20; it flagged \nm{ctl_mucin_d1} mucins. None of the \nm{ctl_flagged_with_cds} flagged
proteins with a retrievable coding sequence showed a DNA-level contradiction: each had a Fickett median of at least
\nm{ctl_flagged_fickett_min} and at most \nm{ctl_flagged_open_max} open alternative frame. We therefore revised the rule so
that D1 counts only together with D2 or D3. A hold-out set fixed before scoring it, all \nm{hold_n} other reviewed proteins of
1,000--4,999 aa with protein-level evidence, held \nm{hold_trigger} proteins with the pattern; the first rule would have
flagged all of them, and none of the \nm{hold_d1_cds} with a retrievable coding sequence met D2 or D3, so the revised rule
flagged none. At the protein level the audit sequences with the pattern overlap the real repetitive proteins (main text Figure~3a).

\paragraph{Can the RNA-seq rule contradict a model?} The ``contradicted'' rule fired for no model in the audit. To see
whether it could, we made deliberately wrong versions of \nm{power_models} well-covered real models (the human titin
control and every eukaryotic model with a ``full'' verdict and at least 100 testable introns): introns shifted by
15--60 bp, or internal exons removed, each setting repeated \nm{power_reps} times at random.
\begin{itemize}[leftmargin=*]
\item With at least 60\,\% of the introns wrong, at least \nm{pw_contr_ge60_min_pct}\,\% of versions were called
contradicted. At 50\,\% wrong, \nm{pw_shift50_pct}\,\% of shifted and \nm{pw_skip50_pct}\,\% of skipped versions were.
\item With 20--40\,\% of the introns wrong, the original rule order called at least \nm{pw_v1_partial_2040_min_pct}\,\% of
versions ``partial'', which counts as a line of support.
\item Under the revised order (the test for unsupported introns with covered flanks first), every version with 30--40\,\%
wrong introns is downgraded. At 20\,\% wrong, \nm{pw_v2_shift20_pct}\,\% of shifted but only \nm{pw_v2_skip20_pct}\,\% of
skipped versions are; at 10\,\%, \nm{pw_v2_10_down_pct}\,\% are.
\item The revision changes no verdict in the real data. The rule still misses wrong introns at a 10\,\% share and most
versions with 20\,\% of exons skipped. Removing every internal exon leaves one intron, which cannot be tested.
\end{itemize}

\paragraph{Known misassemblies and the source-read check (Table~\ref{tab:misassembly}).} From the finished \spn{Nocardia}
sp.\ NBC\_01730 genome and its own reads (all \nm{dna_nbc01730_reads} Illumina pairs and \nm{dna_nbc01730_ont_reads}
nanopore reads) we built three misassemblies whose truth is known (a 10-kb deletion join, a two-copy 4.55-kb repeat collapsed
to one copy, a 5-kb segment duplicated in tandem) and used two correct regions with long real repeats. The first version of
the check (v3) caught the deletion, passed the collapsed repeat as ``supported by the reads'', and left the tandem
duplication at ``supported outside 9.37 kb of repeat'': the repeat exemption hid a real error, though it did not raise a
grade. We then added a copy-number test (v3.1) that fires on at least three consecutive 1-kb windows whose total depth is
above $\sqrt{2}$ or below $1/\sqrt{2}$ of the gene's depth. It caught the collapse and the duplication and left the correct
regions unchanged. Re-judged under v3.1, none of the five real checks changed verdict; over the repeat copies of the checked
genes, depth was \nm{dnactl_repeat_depth_ratio_range} of the gene's depth (\nm{dnactl_omni_repeat_depth_ratio_range} for the
Omnitrophota copies), as expected for correctly assembled copies. The test cannot see a copy-number error shorter than 3 kb.

\begin{table}[ht]
\centering\scriptsize
\caption{The source-read check on known misassemblies and correct repeat regions of the finished \spn{Nocardia} sp.\
NBC\_01730 genome, checked with the genome's own reads. v3.1 adds the copy-number test.}\label{tab:misassembly}
\begin{tabularx}{\textwidth}{L{2.6cm} l Y Y Y}
\toprule
Region & Truth & v3, Illumina & v3, Illumina + nanopore & v3.1, Illumina + nanopore \\
\midrule
\tablebody{tables/tableS2_misassembly.tex}
\bottomrule
\end{tabularx}
\end{table}

\paragraph{Threshold sensitivity (Table~\ref{tab:sensitivity}).} We varied each threshold of Stages 1 and 2, and each rule
revision, one at a time (\nm{sens_variants} settings), and each threshold of the source-read check (\nm{dna_sens_settings}
settings), rerunning the full pipeline in a copy of the repository; the published setting reproduced the released grades
exactly. The six sequences longer than \nm{pkz_len} aa at grade A or B kept their grades in \nm{sens_big_ab_stable} of the
\nm{sens_variants} Stage 1--2 settings. The exceptions: counting a single open alternative frame as D3 would make the two
GC-rich \spn{Nocardia} ORFs and the \spn{Sycon} model D, and requiring 10 junction reads per intron would move the
\spn{Sycon} model to C. Requiring 10 pairs per boundary would move NBC\_00881 from B to C, partly because its long-read
crossing cannot be recounted from the stored result. Counting twin records as independent changes \nm{twin_fix_big} grades above
\nm{pkz_len} aa and \nm{twin_fix_n} below it. Not stable: the split between D and C above \nm{pkz_len} aa, which depends
on the repeat rule (\nm{fin_big_D} D under the revised rule, \nm{sens_d1v1_big_D} under the first), and the number of A
grades among all \nm{n_retained} (\nm{sens_allA_range}), which depends on the RNA thresholds almost entirely through
vertebrate titin isoforms.

\begin{table}[ht]
\centering\scriptsize
\caption{Threshold sensitivity. Each threshold varied alone; ranges include the published value. Counts are sequences
(all \nm{n_retained}; the \nm{n_retained_gt_pkz} longer than \nm{pkz_len} aa). Last column: most sequences changing grade
(of which above \nm{pkz_len} aa).}\label{tab:sensitivity}
\setlength{\tabcolsep}{3pt}
\resizebox{\textwidth}{!}{%
\begin{tabular}{L{3.3cm} L{1.9cm} L{1.7cm} l l L{2.3cm}}
\toprule
Threshold or rule & Published & Tested & All: A / B / C / D & Above: A / B / C / D & Changing \\
\midrule
\tablebody{tables/tableS3_sensitivity.tex}
\bottomrule
\end{tabular}}
\end{table}

\FloatBarrier
\section*{Supplementary Note 3. Case walk-throughs}
\addcontentsline{toc}{section}{Note 3. Case walk-throughs}

\paragraph{KAI4551101.1, a translated tandem repeat.} The model (\nm{ovis_length} aa) from an \spn{Ovis ammon polii}
$\times$ \spn{O.~aries} hybrid assembly is annotated as \nm{ovis_n_exons} exons over \nm{ovis_genomic_span_bp} bp. The
protein repeats with a period of \nm{ovis_protein_period_aa} aa at \nm{ovis_protein_period_identity} identity, and
\nm{ovis_rep25_frac} of its residues lie in repeated 25-mers; it has no recognisable domain. The coding sequence has a GC
fraction of \nm{ovis_gc_cds} and repeats with a period of \nm{ovis_dna_tandem_lag_bp} bp at \nm{ovis_dna_tandem_identity}
identity; its Fickett coding score is \nm{ovis_fickett_median}, in the non-coding range. Two of the five alternative frames
are nearly as open as the annotated one (\nm{ovis_min_stops_per_kcodon_other_frames} stop codons per 1,000 codons, against
about 45 in the reverse frames). One RNA-seq run of \nm{ovis_rna_reads} reads gave \nm{ovis_rna_candidates} candidate reads
on the locus. All \nm{ovis_entries} retained entries from this hybrid show the 22-aa period or a multiple of it and remain D
under the revised rule because their coding sequences fail D2 or D3.

\paragraph{Fish titin pair.} In KAK2904762.1 (\nm{cirrhinus_len} aa, \nm{cirrhinus_exons} exons), \nm{cirrhinus_runs}
RNA-seq runs gave \nm{cirrhinus_reads_unique} uniquely mapped reads. TSR51413.1 (\nm{bagarius_len} aa) restarts at a titin
N-terminus at aa \nm{bagarius_second_copy_start_aa} (first part \nm{bagarius_first_ttna} windows \gene{ttna}-like, second
part \nm{bagarius_second_ttnb} \gene{ttnb}-like); its \nm{bagarius_runs} runs were too shallow to test the join.

\paragraph{The two \spn{Nocardia} chromosomes.} NBC\_01730 (WP\_327112628.1): all \nm{dna_nbc01730_reads} Illumina reads
used; median unique depth over the gene \nm{dna_nbc01730_depth} (gene-to-contig ratio \nm{dna_nbc01730_ratio}); at least
\nm{dna_nbc01730_span_min_testable} uniquely placed pairs span every testable 1-kb boundary; a 9.0-kb repeat at the start of
the gene is spanned by \nm{dna_nbc01730_ont_span_blocks} nanopore reads. NBC\_00881 (WP\_406270362.1): all
\nm{dna_nbc00881_reads} Illumina reads used; at least \nm{dna_nbc00881_span_min_testable} pairs span every testable boundary;
two short stretches that pairs cannot place (1.1 kb together) are spanned by \nm{dna_nbc00881_ont_span_blocks} nanopore
reads. Both are supported by their reads. With an orthologue of consistent length in \spn{N.~arseniciresistens}
(\nm{nbc01730_best_len} aa), NBC\_01730 reaches A; NBC\_00881 reaches B, the same homologue giving it a length ratio of
\nm{nbc00881_ratio}, outside the $\pm$20\,\% rule.

\paragraph{The \nm{orf_len_aa}-aa Omnitrophota ORF.} The three proteins (DAOTDT01, DAOPMG01, DAOOUB01) are
\nm{omni_aa_identity_min_pct}--\nm{omni_aa_identity_max_pct}\,\% identical, on contigs of \nm{omni_contig_lengths} bp; the
coding sequence in DAOTDT01 contains \nm{omni_n_in_cds} N. All reads of each sample were used (\nm{dna_tdt_reads},
\nm{dna_pmg_reads} and \nm{dna_oub_reads}). Over the gene, median unique depth was \nm{dna_tdt_depth}, \nm{dna_pmg_depth}
and \nm{dna_oub_depth}; the gene-to-contig depth ratio \nm{dna_tdt_ratio}, \nm{dna_pmg_ratio} and \nm{dna_oub_ratio};
discordant pairs at most \nm{dna_tdt_discordant_pct}\,\% of any window; no clip pile-up broke read continuity; no consensus
change breaks the frame. The two repeat copies (\nm{rep_copy_bp} bp each) lie \nm{rep_lag_bp} bp apart (an in-frame
distance of \nm{rep_lag_codons} codons) and differ at \nm{rep_diffs} positions, the same coding positions in all three
assemblies. The first version of the check reported a ``break'' with zero spanning pairs at one 1-kb boundary in two
assemblies; each boundary lay in the middle of a repeat copy, where fragments of the library's size (median insert
\nm{dna_oub_insert}--\nm{dna_tdt_insert} bp) cannot be placed uniquely, and the 1-kb window mask could not see a repeat
shorter than a window. The revised check simulates the library on the assembly and asks how many unique spanning pairs a
correct assembly would give at each boundary; at all three flagged boundaries the expected number was 0, and 0 was observed.
DAOTDT01 and DAOPMG01 are ``supported outside 1.3 kb of repeat''; DAOOUB01 is ``inconclusive'' because one boundary at the
edge of a copy has \nm{dna_oub_weakest_obs} pairs where about \nm{dna_oub_weakest_exp} are expected (the minimum is 5). If the
10.4-kb segment between the copies were placed elsewhere, the frame would remain open at \nm{orf_len_if_excised} aa.

\section*{Supplementary Note 4. RNA-seq stage in detail}
\addcontentsline{toc}{section}{Note 4. RNA-seq stage in detail}

\paragraph{Runs.} \nm{rna_runs_kept} run results from \nm{rna_species} species were kept (\nm{rna_runs_v2} from the final
pipeline, which read \nm{rna_reads_billion}$\times10^9$ reads, and \nm{rna_runs_v1} from a first version); re-run results
replaced earlier results of the same run. \nm{rna_runs_stream_error} runs ended because the archive closed the connection;
they contributed the reads received before the drop (a median of \nm{rna_stream_error_frac_median_pct}\,\% of each run, as
little as \nm{rna_stream_error_frac_min_pct}\,\%). Independently, \nm{rna_runs_saturation} runs ended at the saturation
rule and \nm{rna_runs_readcap} at the cap of 60 million reads per file. Junction counts are therefore lower bounds; partial
reading cannot create support. The \nm{rna_runs_v1} first-version results could have counted some reads twice after a
restarted download; this affects read counts for PKZILLA-1, not the presence of its junctions.

\paragraph{Controls and verdicts.} Human titin: \nm{titin_introns_counted_supported} of \nm{titin_n_introns} introns
supported (\nm{titin_introns_unique_support} by unique reads), none of its \nm{titin_testable_introns} testable introns
unsupported. PKZILLA-1: all \nm{pkz_n_introns} supported, \nm{pkz_introns_unique_support} by unique reads and
\nm{pkz_introns_multi_only}, inside internally repeated modules, only by multi-mapped reads. Verdicts: \nm{rna_full} full,
\nm{rna_partial} partial, \nm{rna_weak} weak, \nm{rna_noexpr} no expression; \nm{rna_noruns} models had no usable run and
\nm{rna_nottested} could not be tested; none was contradicted. Four titin models first looked weakly supported because
their runs had been cut by the saturation rule; with whole runs the \spn{Microcaecilia}, \spn{Rana} and \spn{Solea} models
are ``partial'', with 1 unsupported intron among \nm{microcaecilia_testable}, \nm{rana_testable} and \nm{solea_testable}
testable introns, and the \spn{Pantherophis} model remains limited by coverage (0 of \nm{pantherophis_testable}). One
model placed on a different haplotype assembly had \nm{liftover_moved} of its \nm{liftover_n_introns} introns shifted by
a few bases and a false ``contradicted'' verdict; a whole-locus lift-over corrected it, and a guard now flags such shifts.

\paragraph{The repeat-flag models.} The \nm{rep_models} models graded D under the first repeat rule were never sent to the
RNA stage. After the rule was revised they were tested in a separate run (\file{cloud/bundle_rep_rna}), under rules written
down before the run (\file{cloud/bundle_rep_rna/RULES.txt}): the published rules unchanged, with multi-mapped reads not
counted (repeat flag); ``full'' or ``partial'' moves C to B; ``contradicted'' moves C to D only if it also holds on the
introns whose junction sequence (50 bp either side) is unique within the model, or through junctions leaving the gene; all
other verdicts leave C. \nm{rep_testable} models in \nm{rep_jobs} species jobs could be tested (\nm{rep_runs} runs,
\nm{rep_gb} GB, downloaded to disk with resumed retries); \nm{rep_untestable_noruns} have no Illumina RNA-seq run in ENA
(\nm{rep_noruns_plectranthus} \spn{Plectranthus ornatus}, \nm{rep_noruns_merlangius} \spn{Merlangius merlangus}, \nm{rep_noruns_coilia} \spn{Coilia grayii}) and \nm{rep_untestable_nolocus}
(\spn{Fasciola hepatica} UniParc entries) have no gene model with coordinates. All \nm{rep_run_results} run results came
back, none skipped. Before the run, we counted for each model the introns whose junction is unique within the model
(Table~\ref{tab:reprna}): all for \spn{Arctogadus} and \spn{Owenia}, \nm{mapp_schisto1_pct}\,\% for \spn{Schistosoma}, and
\nm{mapp_syn1_pct}\,\% or less for the \spn{Synaphobranchus}, \spn{Acnodon} and stickleback models. For the latter,
``partial'' was out of reach, and the stickleback and \spn{Acnodon} models could not reach 10 testable introns at all. The
power test calibrated the rules on non-repetitive models; for repeat models an error confined to the repeated copies,
including a wrong number of copies, is invisible, and a clean result means only that the uniquely mappable parts are
spliced as annotated.

Results (Table~\ref{tab:reprna}): both \spn{Schistosoma haematobium} isoform models are ``partial''
(\nm{rep_schisto1_supported} of \nm{rep_schisto1_introns} and \nm{rep_schisto2_introns} introns supported by unique reads,
none unsupported among \nm{rep_schisto1_testable} testable introns) and move from C to B; \spn{Owenia} (\nm{rep_owenia_supported}
of \nm{rep_owenia_introns}) and \spn{Arctogadus} (\nm{rep_arcto_supported} of \nm{rep_arcto_introns}) are ``weak''; the six
others had at most \nm{rep_noexpr_max_unique} uniquely placed reads (``no-expression''). No model was contradicted, so the
guard was never applied. \emph{Disclosure:} the driver's download mode was tested on a Mac with one small \spn{Schistosoma}
run (ERR11178326) before \file{RULES.txt} was committed; its result (5 introns supported in that run) was seen after the
guard code had been written (file times 11:02--11:03 UTC; result 11:03:37 UTC) and was not used. The verdict does not
depend on that run: another run alone (ERR10016938) supports \nm{rep_schisto_err10016938} introns.

\begin{table}[ht]
\centering\scriptsize
\caption{RNA-seq of the testable repeat-flag models. Unique junctions: introns whose junction sequence occurs once in the
model (computed before the run). Supported: introns with $\geq$ 3 exact uniquely placed reads. Testable: both flanking exons
at unique depth $\geq$ 10 in one run (unsupported among them after the slash).}\label{tab:reprna}
\resizebox{\textwidth}{!}{%
\begin{tabular}{l l r r l r r r l l l}
\toprule
Model & Organism & aa & Introns & Unique junctions & Runs & Unique reads & Supported & Testable & Verdict & Grade \\
\midrule
\tablebody{tables/tableS4_repeat_rna.tex}
\bottomrule
\end{tabular}}
\end{table}

\FloatBarrier
\section*{Supplementary Note 5. Proteomics}
\addcontentsline{toc}{section}{Note 5. Proteomics}

\paragraph{PXD025008, deposited peak lists.} We first searched the six deposited MudPIT peak lists (\nm{prot_spectra}
spectra from \nm{prot_fractions} fractions) with Sage against the authors' database, with the organism's earlier genome-bin
proteins replaced by the \nm{prot_closed_genome_proteins} proteins of the closed genome (CP019384.1), QAT16263.1 replaced by
its RefSeq version WP\_128698900.1 (QAT16263.1 is the same sequence started 9 codons later), and a reversed and a shuffled
copy of the giant added (\nm{prot_db_entries} entries). At 1\,\% FDR: \nm{prot_psms} PSMs, \nm{prot_peptides} peptides and
\nm{prot_proteins} proteins; for the giant \nm{prot_giant_peptides} unique peptides (\nm{prot_giant_psms} PSMs, protein
$q = $ \nm{prot_giant_protein_q}) covering \nm{prot_giant_residues} residues (\nm{prot_giant_cov_pct}\,\%) in
\nm{prot_giant_regions} regions; per length fifth \nm{prot_giant_fifths}; the last at residue \nm{prot_giant_last_end}.
Controls: reversed and shuffled copies \nm{prot_rev_psms} PSMs (best protein $q$ \nm{prot_rev_q} and \nm{prot_shuf_q});
elongation factor Tu (\nm{prot_eftu_len} aa) \nm{prot_eftu_peptides} unique peptides covering \nm{prot_eftu_cov_pct}\,\%,
ribosomal protein L7/L12 (\nm{prot_l7_len} aa) \nm{prot_l7_peptides} covering \nm{prot_l7_cov_pct}\,\%, both chosen by name
before their results were seen. Among the \nm{prot_top_n} \spn{Velamenicoccus} proteins with the most PSMs, the giant ranks
\nm{prot_top_rank_giant}th; the other proteins are covered to \nm{prot_top_cov_range}\,\% by unique peptides. In one of
the authors' six identification files (Scaffold export), the \nm{prot_authors_pass_peptides} peptides on the giant that
pass their threshold (their database holds a \nm{prot_authors_len}-aa version) all start in its first fifth.

\paragraph{The deposited peak lists are a selection.} The peak lists hold \nm{prot_scan_share_pct}\,\% of the scan numbers
of their fractions, and \nm{prot_id_rate_pct}\,\% of their spectra were identified at 1\,\% FDR. The raw files settle it:
they hold \nm{raw_ms2} MS/MS spectra (\nm{raw_instrument}), of which the peak lists contain \nm{raw_mgf_share_pct}\,\%.

\paragraph{PXD025008, raw files.} All \nm{raw_files} RAW files (\nm{prot_raw_gb} GB, SHA-1 checked) were converted
without error and searched with the same database and settings (\nm{raw_wall_min} min on the laptop). At 1\,\% FDR:
\nm{raw_psms} PSMs, \nm{raw_peptides} peptides, \nm{raw_proteins} proteins. Giant: \nm{raw_giant_peptides} unique peptides
(\nm{raw_giant_psms} PSMs, protein $q = $ \nm{raw_giant_protein_q}), \nm{raw_giant_residues} residues
(\nm{raw_giant_cov_pct}\,\%), per length fifth \nm{raw_giant_fifths}; \nm{raw_giant_first6000} start in the first 6,000 aa;
the others start at \nm{raw_giant_after6000}; none ends after residue \nm{raw_giant_last_end}. \nm{raw_both} peptides are
shared with the peak-list search, \nm{raw_only_raw} are found only in the raw files and \nm{raw_only_mgf} only in the peak
lists, among them the one at \nm{raw_mgf_only_after6000} (peptide $q = $ \nm{prot_peak_last_q} in the peak-list search).
Controls: reversed and shuffled copies \nm{raw_rev_psms} PSMs (best protein $q$ \nm{raw_rev_q} and \nm{raw_shuf_q});
elongation factor Tu \nm{raw_eftu_peptides} unique peptides (\nm{raw_eftu_cov_pct}\,\%), L7/L12 \nm{raw_l7_peptides}
(\nm{raw_l7_cov_pct}\,\%). Conversion used ThermoRawFileParser v.2.0.0-dev for macOS arm64, unpacked in the project folder;
macOS refused to run the unsigned binary, so it was code-signed locally (ad hoc), which does not change what it computes.

\paragraph{\spn{Sycon ciliatum}.} PRIDE lists two projects for the species. PXD031991 \citep{graubove2022} holds histones
acid-extracted from nuclei and derivatised (propionylation) before digestion; it is not suited to detecting a giant protein
and was not searched. PXD060105 \citep{voigt2025} is a proteome of the spicule organic matrix (tissue removed with sodium
hypochlorite, matrix decalcified, proteins separated by gel). The authors' database (\nm{sy_db_authors} transcript-derived
proteins) contains the model XP\_065194412.1 as \nm{sy_fragments} partial fragments, the longest \nm{sy_fragment_max_len}
aa; their report of \nm{sy_authors_n} proteins does not include it, and their paper does not mention it. We replaced the
fragments by the full model, added a reversed and a shuffled copy, and searched the \nm{sy_raw_files} RAW files
(\nm{sy_ms2} MS/MS spectra, Orbitrap). The dataset is small: the lowest peptide- and protein-level $q$ of any target were
\nm{sy_min_pep_q} and \nm{sy_min_prot_q}, so nothing can pass 1\,\% at those levels. At 1\,\% spectrum-level FDR the search
gave \nm{sy_psms_spec} PSMs and \nm{sy_peptides_spec} peptides and recovered \nm{sy_authors_found} of the authors'
\nm{sy_authors_n} proteins (their top protein, Calcarin 3, \nm{sy_calc3_psms} PSMs; tubulin beta \nm{sy_tub_psms}); the model
obtained \nm{sy_giant_psms_spec} PSMs (\nm{sy_giant_psms_any} at any $q$, against \nm{sy_rev_psms_any} and
\nm{sy_shuf_psms_any} for its reversed and shuffled copies). The absence is uninformative for a protein that is not expected
in this fraction.

\begin{table}[ht]
\centering\scriptsize
\caption{Proteomics searches. PSMs and peptides at 1\,\% FDR (spectrum, peptide and protein level); $^{a}$PXD060105 at
spectrum level only, because the dataset cannot reach 1\,\% at the other levels. Last residue: end of the last unique
peptide on the giant (share of its length). Controls: PSMs of the reversed / shuffled giant at 1\,\% FDR.}\label{tab:prot}
\resizebox{\textwidth}{!}{%
\begin{tabular}{l l r r r r r l l}
\toprule
Search & Giant & MS/MS spectra & PSMs & Peptides & Giant peptides & Giant PSMs & Last residue & Controls \\
\midrule
\tablebody{tables/tableS5_proteomics.tex}
\bottomrule
\end{tabular}}
\end{table}

\FloatBarrier
\section*{Supplementary Note 6. Limitations in full}
\addcontentsline{toc}{section}{Note 6. Limitations in full}
\begin{enumerate}[leftmargin=*]
\item \textbf{The rules are ours and were revised during the work} (Note~1). Thresholds were set by us on two controls, not
by an outside annotator, and should be reviewed.
\item \textbf{Partial downloads.} \nm{rna_runs_stream_error} RNA-seq run results ended because the archive closed the
connection; junction counts are lower bounds, and a missing junction in a partly read run is weak evidence, which is why
only introns with covered flanks count against a model (Note~4).
\item \textbf{Tissue and coverage.} RNA-seq was taken from whatever tissues ENA offered; several models were tested only in
tissues that probably do not express them (for example titins in blood or gonad). A ``weak'' or ``no-expression'' verdict
is not evidence against a model.
\item \textbf{Little protein evidence.} Two proteomic datasets were searched and a third examined. For PXD025008 the raw
files were searched; our search limits (one missed cleavage, peptides of 7--40 aa) were set to fit a laptop's memory. For
the other entries PRIDE was searched only by species name \citep{perezriverol2025}. RNA-seq shows that introns are spliced
as annotated, and peptides that parts of a protein are made; neither shows that the full-length protein is made as one
chain.
\item \textbf{Repeats limit what short reads can test.} For the Omnitrophota ORF, one 1.2-kb repeat cannot be crossed by
the available libraries. Junction counting inside repeated modules needs multi-mapped reads (PKZILLA-1), so models flagged
as repeats do not get credit for such reads, and most of their introns cannot be tested at all (Note~4).
\item \textbf{Locus placement.} RNA-stage loci were placed on the downloaded assembly by a sequence probe and a constant
offset; across haplotypes this shifted \nm{liftover_moved} of one model's \nm{liftover_n_introns} introns and produced a
false ``contradicted'' verdict, corrected by a whole-locus lift-over. A guard now flags any model whose reads show its
junctions a few bases away; no other model showed this.
\item \textbf{What the controls do not cover.} The repeat-rule controls used reviewed UniProtKB proteins, which favour
well-studied organisms. The misassembly controls used one bacterial genome with deep coverage. The copy-number test cannot
see errors shorter than 3 kb, so it could not detect a collapse of the 1.2-kb Omnitrophota repeat; the measured depth over
those copies is what two correct copies would give, which is consistent with, not proof of, the assembled arrangement.
\item \textbf{Scope.} MGnify and other metagenomic protein catalogues were not included. Of the bacterial models, only the
three discussed were checked against their source reads. Database contents change, so every count refers to the stated
snapshot dates. Orthology was judged by length consistency of DIAMOND hits in UniProtKB, which favours well-sampled clades.
\item \textbf{Independence.} The three Omnitrophota assemblies come from three samples but one study and one assembly
pipeline; agreement between them does not exclude a systematic error at the repeat.
\end{enumerate}

\section*{Supplementary Methods}
\addcontentsline{toc}{section}{Supplementary Methods}

\paragraph{Inventory.} UniProtKB (release 2026\_04) and UniParc were queried through the UniProt REST API for sequences
longer than \nm{titin_len} aa \citep{uniprot2025}; UniParc cross-references were fetched per entry. NCBI Protein was queried
with the Entrez length filter \texttt{35992:100000000[SLEN]} on 11 November 2026 \citep{sayers2026,oleary2016}; GenPept
records were downloaded for coding-sequence coordinates, gene identifiers and RefSeq model evidence. Sequences were merged
by MD5 checksum; withdrawn or patent/synthetic records and titles marked ``partial'' were excluded; loci were defined by
shared NCBI GeneID or overlapping genomic coordinates.

\paragraph{Stage 0: sequence flags.} For each protein: the share of X residues and of low-complexity 12-aa windows; the
share of residues inside 25-mers that occur at least twice (rep25); the best tandem period and the self-identity at that
lag; coverage by a panel of \nm{hmm_panel} Pfam profile HMMs \citep{mistry2021} (polyketide and peptide synthases, Ig/FN3
and titin domains, adhesin repeats and other families typical of giant proteins), retrieved from the InterPro API on
11 November 2026 and searched with PyHMMER \citep{eddy2011,larralde2023} at gathering thresholds.

\paragraph{Stage 1: locus checks and grades.} Locus DNA was fetched from INSDC/RefSeq. For each model we measured the
number and size of exons, contig length and distance to contig ends, N runs in the coding sequence, GC by codon position,
the Fickett TESTCODE statistic in 300-bp windows \citep{fickett1982}, stop codons per 1,000 codons in the five alternative
frames and the DNA tandem period. Each protein was cut into 1,000-aa windows and searched with DIAMOND 2.2.8
(\texttt{--fast}, $E \leq 10^{-10}$) \citep{buchfink2021} against UniProtKB Swiss-Prot plus unreviewed entries of at least
3,000 aa; we recorded the best homologue in another species and the ratio of lengths. A fusion test asked whether disjoint
segments of the model each cover at least 70\,\% of a different complete reference protein.
\begin{itemize}[leftmargin=*]
\item \textbf{D1:} self-identity at the best period $\geq$ 0.6, or rep25 $\geq$ 0.8, with domain coverage $<$ 5\,\%; a
contradiction only together with D2 or D3, otherwise recorded as a repeat flag (first version: the pattern alone).
\item \textbf{D2:} $\geq$ 50\,\% of coding windows non-coding by Fickett, with domain coverage $<$ 5\,\%.
\item \textbf{D3:} $\geq$ 2 of the 5 alternative frames with $<$ 5 stops per 1,000 codons, with domain coverage $<$ 5\,\%.
\item \textbf{D4:} $\geq$ 1\,\% N in the coding sequence, or a model joined across contigs.
\item \textbf{D5:} $\geq$ 2 disjoint segments, each covering $\geq$ 70\,\% of a different complete reference protein of
$\geq$ 1,000 aa and at most 60\,\% of the query length, with no homologue in another species within $\pm$20\,\% of the query
length; polyketide and nonribosomal peptide synthases are exempt.
\item \textbf{Structure OK:} no D rule; no N in the coding sequence; the gene at least 1 kb from both contig ends; a coding
signal (Fickett median $\geq$ 0.95 with $\leq$ 20\,\% of windows at or below 0.74; for GC-rich bacterial genomes, GC2 as the
lowest codon position).
\item \textbf{S1:} a homologue in another species hit on $\geq$ 50\,\% of the windows, with length within $\pm$20\,\%.
\item \textbf{S2:} a near-identical copy ($\geq$ 95\,\% identity on $\geq$ 90\,\% of windows, length within $\pm$5\,\%) in a
different assembly. S1 and S2 do not count a record's twin from the same genome: species names are compared without
``Candidatus'' and strain text, and a RefSeq contig (NZ\_) belongs to the same assembly as its INSDC source.
\item \textbf{S3:} NCBI's own report that RNA-seq supports all introns. \textbf{S4:} a curated RefSeq record.
\item \textbf{S6:} at least 3 unique peptides at 1\,\% FDR (spectrum, peptide and protein level), with at least one in each
fifth of the model's length; peptides in fewer fifths are noted but are not a support line. (S5 labels Stage 2 transcript
support in the data files.)
\item \textbf{Grades:} A = structure OK, plus S3 or S4, plus S1 or S2; B = structure OK plus any support line; C = no
contradiction and no support line; D = any contradiction.
\end{itemize}
Human titin NP\_001254479.2 was graded as a control and is not counted among the \nm{n_retained}. Stage 1 used about
\nm{stage1_cpu_hours} CPU-hours on a laptop.

\paragraph{Stage 2: RNA-seq junction evidence.} For eukaryotic models at grade B or C, the fusion checks and the controls,
we selected four public Illumina RNA-seq runs per species (six for the controls and the largest models) from ENA
\citep{yuan2026}: 0.2--6~GB each, not single-cell, spread across studies, muscle first for titin-like models. Genome
assemblies were downloaded from NCBI, or from ENA for assemblies NCBI had not yet linked. Runs were streamed as FASTQ over
HTTPS (downloaded to disk with resumed retries for the repeat-flag stage), pre-filtered with minimap2 2.31 (\texttt{-x sr})
\citep{li2018} against the coding exons $\pm$100 bp, and the candidate reads aligned with \texttt{-x splice:sr} to the whole
genome, which serves as a decoy for paralogous and repeated sequence; alignments were handled with samtools
\citep{danecek2021}. For each annotated intron we counted spliced reads that match it exactly, separately for uniquely
placed (MAPQ $\geq$ 10) and multi-mapped reads. An intron is supported with $\geq$ 3 exact reads; multi-mapped reads count
only for models without a repeat flag; an intron is testable if both flanking exons reach a mean unique depth of 10 in at
least one run. Verdicts in order: no-expression ($<$ 50 unique reads); contradicted ($\geq$ 10 testable introns with
$\geq$ 50\,\% unsupported, or more reads on junctions that leave the gene than on the annotated junctions); under-supported
where covered ($\geq$ 10 testable introns with $\geq$ 20\,\% unsupported; tested after full and partial in the first
version); full ($\geq$ 90\,\% supported); partial ($\geq$ 50\,\%); otherwise weak. Full raises C to B and B to A where
Stage 1 had S1 or S2; partial raises C to B; under-supported lowers A to B and B to C; contradicted gives D (for the
repeat-flag stage only with the guard of Note~4). Each run stopped at the end of the data, at 60 million reads per file,
after 1,800 s of download (stream mode), or, in the first pass only, at a saturation level of candidate reads; runs cut short
by the saturation rule were repeated without it. Loci placed on an assembly other than the model's own were checked by
aligning the source locus to the target locus, and junction reads were counted at lifted coordinates where they moved.

\paragraph{Stage 3: source-read check (rules v3 and v3.1).} Illumina pairs (and nanopore reads, where available) from the
assembly's BioSample were streamed from ENA, or read from downloaded copies after repeated stream failures, and mapped with
minimap2 to the whole assembly; only primary alignments with $\leq$ 3\,\% mismatches (Illumina) or $\leq$ 15\,\% (nanopore)
were kept, and a mapping counted only if the number of reads processed reached ENA's read count. Measures in 1-kb windows
over the gene: unique (MAPQ $\geq$ 20) and all-alignment depth; uniquely placed proper pairs (fragment $\leq$ 3$\times$ median
insert) spanning each window boundary by $\geq$ 20 bp on both sides; clip pile-ups and whether physical coverage is lost
there; discordant pairs; gene-to-contig depth ratio; positions where $\geq$ 50\,\% of reads carry an indel or a stop-creating
substitution. Error-free pairs with the library's median insert and read length were cut from the assembly every 2 bp and
mapped back; the expected number of unique spanning pairs at a boundary is the share of simulated pairs placed uniquely
across it times the gene's median observed spanning count. A boundary is tested only if this expectation is $\geq$ 5;
stretches below it must be crossed as a whole by $\geq$ 5 pairs or $\geq$ 3 nanopore reads with 500-bp flanks, otherwise they
are reported as unresolved repeats. Supported: all testable boundaries spanned by $\geq$ 5 pairs ($\geq$ 3 long reads),
$\geq$ 99\,\% of single-copy bases at depth $\geq$ 5, no clip pile-up with lost continuity, discordant pairs $\leq$ 10\,\%,
depth ratio 0.5--2, no frame-breaking change. Against: a testable boundary with 0 pairs, a clip pile-up with lost
continuity, $>$ 10\,\% discordant pairs, an out-of-range depth ratio, a frame-breaking change, or (since v3.1) $\geq$ 3
consecutive windows whose all-alignment depth is above $\sqrt{2}$ or below $1/\sqrt{2}$ of the median over the gene's
single-copy windows. Supported outside N kb of repeat: supported except across unresolved repeats; anything else
inconclusive. Every check ran a negative control (an inverted join at the gene midpoint). Supported with S1 or S2 gives A,
supported alone B, against D.

\paragraph{Proteomics.} PXD025008 \citep{kizina2022}: peak lists, raw files and the authors' database were downloaded
from PRIDE \citep{perezriverol2025} (\file{proteomics/run_download.sh}, \file{run_download_raw.sh}); the database was rebuilt
by \file{proteomics/build_db.py}. Sage 0.14.7 \citep{lazear2023} (release binary for macOS arm64, which reports itself as
0.14.6; 4 threads): trypsin (K or R, not before P), at most 1 missed cleavage, peptides of 7--40 aa, precursor tolerance
10 ppm, fragment tolerance 0.5 Da, isotope errors 0--1, variable methionine oxidation (at most 2), no fixed modification,
reversed decoys generated by Sage. Raw files were converted to indexed mzML by ThermoRawFileParser \citep{hulstaert2020}
with vendor peak picking (\file{proteomics/run_search_raw.sh}). PXD060105 \citep{voigt2025}: the same pipeline
(\file{proteomics/sycon/}), with fixed carbamidomethylation of cysteine and a fragment tolerance of 20 ppm for the Orbitrap
spectra. Identifications were kept at spectrum, peptide and protein $q \leq 0.01$ (Sage target-decoy). A peptide counts as
unique if its sequence, with I and L treated as equal, occurs in no other forward entry of the database. Positions,
coverage and per-fifth counts are computed from peptide start positions (\file{analyze_verified.py},
\file{verify_extra.py}, \file{verify_raw.py}, \file{sycon/analyze.py}).

\paragraph{Controls of the grading rules.} The repeat-rule design and hold-out sets were defined by UniProtKB queries fixed
before scoring (\file{controls/d1/}): reviewed entries with evidence at protein level; silk proteins (fibroin, spidroin or silk
in the name, at least 300 aa, excluding silk gland factors) plus unreviewed complete spidroins of at least 2,500 aa with
protein-level evidence; mucins of at least 2,000 aa; bacterial adhesins of at least 2,000 aa; all proteins of at least
5,000 aa; the hold-out set all other reviewed entries of 1,000--4,999 aa. Scoring used the Stage 1 code paths; coding
sequences came from ENA (at most three EMBL CDS records per entry, translation $\geq$ 95\,\% identical to the protein). The
RNA-seq power test (\file{controls/rna_power/}) scored wrong models against the stored junction counts of the same runs
(annotated junctions plus each run's 50 most frequent other junctions), with testability taken from the original flanking
exons; fixed random seed. Known misassemblies (\file{controls/dna/}) were edited copies of GCA\_035920445.1 checked with its
own Illumina (SRR21428551) and nanopore (SRR21563524) reads. Threshold sensitivity (\file{controls/sensitivity/}) reran the
full pipeline in a throw-away copy of the repository, once per setting; source-read thresholds were varied by re-judging the
stored results.

\paragraph{Leaderboard.} One row per locus, the longest live model representing it. Copies in other assemblies (S2), and
records of one organism sharing $\geq$ 90\,\% of the shorter protein's 10-mers (haplotypes, GenBank and RefSeq versions,
isoforms; grades A--C only), are merged and listed; merged records are not counted as independent support.

\paragraph{Software.} Python 3 (standard library), minimap2 2.31, samtools 1.25, DIAMOND 2.2.8, PyHMMER 0.12.3, Sage
0.14.7, ThermoRawFileParser v.2.0.0-dev, matplotlib; LaTeX built with tectonic.

\section*{Supplementary Table S1}
\addcontentsline{toc}{section}{Supplementary Table S1}
The leaderboard (\file{supplementary/Table_S1_leaderboard.tsv}): \nm{lb_loci} graded loci and \nm{lb_excluded_rows}
excluded sequences, one row each. Column dictionary:

{\small
\begin{longtable}{L{5.7cm} p{9.6cm}}
\toprule
Column & Description \\
\midrule
\endhead
\tablebody{tables/tableS1_columns.tex}
\bottomrule
\end{longtable}}

\bibliographystyle{unsrtnat}
\bibliography{refs}
\end{document}
